% !TEX root = ../algebra_02.tex
\section{Орбиты и стабилизаторы}

\begin{Definition}{Орбита}{}
	Пусть $G \curvearrowright M$.
	\textbf{Орбитой} элемента $m \in M$ называется множество:
	\[ Gm = \{gm \mid g \in G\} \] 
\end{Definition}

\begin{Lemma}{О непересечении орбит}{}
	Для любых $m_1, m_2 \in M$ либо $Gm_1 = Gm_2$, либо $Gm_1 \cap Gm_2 = \emptyset$.
\end{Lemma}

\begin{proof}
	Предположим, что $Gm_1 \cap Gm_2 \neq \emptyset$. Тогда существует $y \in Gm_1 \cap Gm_2$, то есть найдутся $g_1, g_2 \in G$ такие, что $y = g_1 m_1 = g_2 m_2$.
	Отсюда выразим $m_2$:
	\[ m_2 = g_2^{-1}(g_2 m_2) = g_2^{-1}(g_1 m_1) = (g_2^{-1}g_1)m_1 \] 
	Тогда для любого $h \in G$ выполнено:
	\[ h m_2 = (h g_2^{-1}g_1)m_1 \in Gm_1 \] 
	Следовательно, $Gm_2 \subset Gm_1$.
	Аналогично доказывается, что $Gm_1 \subset Gm_2$. Значит, $Gm_1 = Gm_2$.
\end{proof}

\begin{Remark}{Связь с классами эквивалентности}{}
	Орбиты являются классами эквивалентности на $M$ по отношению: $m_2 \sim m_1 \iff \exists g \in G \colon m_2 = gm_1$.
\end{Remark}

\begin{Definition}{Стабилизатор}{}
	Пусть $G \curvearrowright M$ и $m \in M$. \textbf{Стабилизатором} элемента $m$ называется множество:
	\[ St_m = \{g \in G \mid gm = m\} \] 
	Очевидно, что $St_m \le G$ (является подгруппой), так как если $gm=m$, то умножая слева на $g^{-1}$, получаем $m = g^{-1}m$, то есть $g^{-1} \in St_m$.
\end{Definition}

\begin{Proposition}{Теорема об орбите и стабилизаторе}{}
	Если $G \curvearrowright M$, то мощность орбиты равна индексу подгруппы стабилизатора:
	\[ |Gm| = (G : St_m) \]
\end{Proposition}

\begin{proof}
	Построим отображение $\alpha \colon G/St_m \to Gm$, заданное правилом $\alpha(gSt_m) = gm$.
	\begin{itemize}
		\item \textbf{Корректность:} Если $g' = gh$ для некоторого $h \in St_m$, то $g'm = g(hm) = gm$.
		Отображение не зависит от выбора представителя смежного класса.
		\item \textbf{Сюръективность:} Очевидна по определению орбиты.
		\item \textbf{Инъективность:} Пусть $\alpha(g_1 St_m) = \alpha(g_2 St_m)$. Тогда $g_1 m = g_2 m \implies m = g_1^{-1}g_2 m \implies g_1^{-1}g_2 \in St_m \implies g_2 \in g_1 St_m \implies g_1 St_m = g_2 St_m$.
	\end{itemize}
	Построенная биекция доказывает утверждение.
\end{proof}

\begin{Corollary}{Следствие для конечных групп}{}
	Если группа $G$ конечна, то для любого $m \in M$ размер орбиты $|Gm|$ делит порядок группы $|G|$.
\end{Corollary}
